\section{Quantitative Evaluation Framework}
\label{sec:evaluation}

Each objective is paired with a testable hypothesis (Table~\ref{tab:hypotheses}).
Results are grouped into three pillars:
\textbf{P1 AI accuracy} (O3--O5),
\textbf{P2 Semantic fidelity} (O2, O5),
\textbf{P3 Trustworthiness} (O1, O6).

\begin{table*}[t]
  \centering
  \caption{Research hypotheses H1--H6 with baselines and target metrics.}
  \label{tab:hypotheses}
  \small
  \begin{tabular}{clp{3.5cm}p{2.8cm}p{2.5cm}}
    \toprule
    ID & Hypothesis (summary) & Baseline & Primary metrics & Target \\
    \midrule
    H1 & Warehouse lineage $\uparrow$ without high ingest cost & Direct \texttt{cim\_*} ingest & Link rate; overhead & Link $\uparrow$; overhead $\leq Y$\% \\
    H2 & VKG preserves SQL answers & Raw SQL & SPARQL F1; coverage & F1 $\geq Z$ \\
    H3 & CIM benchmark large \& executable & Spider/BIRD subset & \#pairs; executability & $\geq$88K; $\geq$99\% \\
    H4 & Fine-tuned LLM beats general models & GPT-4o-mini; SQLCoder & EX; EA; $\Delta$EX & EX$\geq$85\%; EA$\geq$90\% \\
    H5 & Agents reduce OBDA effort & Manual \texttt{.obda} & Mapping F1; automation \% & F1$\geq$0.7 \\
    H6 & Confidence is informative & Paper~1 single value & $\rho$; stock $\Delta$ & $\rho \leq -0.5$ \\
    \bottomrule
  \end{tabular}
\end{table*}


\subsection{Metrics by objective}

\textbf{O1.} Provenance link rate, metadata completeness, ingest success rate, ingest overhead per 1k features.

\textbf{O2.} Schema coverage (\% mapped columns), SPARQL answer preservation F1, query success rate, p95 latency.

\textbf{O3.} Dataset size, taxonomy balance, gold SQL executability (\%).

\textbf{O4.} EM, EX, Deep EM, SC, EA; $\Delta$EX vs.\ GPT-4o-mini and SQLCoder-7B zero-shot; ablations (architecture, schema prompt, model size).

\textbf{O5.} Mapping F1 vs.\ gold \texttt{.obda}, source SQL EX, SPARQL preservation F1, automation rate (\% mappings accepted without edit).

\textbf{O6.} Method coverage, provenance coverage, agreement rate at tolerance $\tau$ (Table~\ref{tab:tolerances}), median absolute deviation (MAD), conflict rate, Spearman $\rho$ between confidence and cross-method spread, calibration error on truth subset, UBEM stock demand sensitivity.

\begin{table}[t]
  \centering
  \caption{Method agreement tolerances $\tau$ (case study).}
  \label{tab:tolerances}
  \small
  \begin{tabular}{lc}
    \toprule
    Feature & Tolerance $\tau$ \\
    \midrule
    \texttt{building\_height} & $\pm 2$\,m \\
    \texttt{building\_area} & $\pm 10$\% \\
    \texttt{building\_type} & exact match \\
    \texttt{building\_volume} & $\pm 15$\% \\
    \bottomrule
  \end{tabular}
\end{table}


\subsection{Ground-truth subsets for calibration}

Where full city-scale labels are unavailable, we use partial truth subsets:
(1)~buildings with both raster and OSM height;
(2)~a LiDAR or field reference sample ($n \geq 50$);
(3)~metered or EPC-labelled buildings when available.
